大型语言模型的强化学习（RL）对训练与推理引擎之间的微小差异极为敏感，这种现象通常被称为训练-推理不匹配（TIM）。  
然而，完全消除TIM是不切实际的，因为这将带来巨大的回放效率损失。  
在本文中，我们表明，在TIM存在的情况下，RL的不稳定性主要由\textit{漂移}引起：一种训练与推理引擎之间持续存在的偏差，且随着每一步训练不断累积。  
我们推导出一个加性的\textit{“得分中心化”}修正项，该修正项通过抵消漂移来稳定TIM下的强化学习。  
在从0.6B到30B参数的模型训练中，仅使用得分中心化的方法在量化条件下即可达到或超越基于重要性采样的方法，且当不匹配程度加剧时，性能差距进一步扩大。  
由于该修正为加性形式，得分中心化可与重要性采样相结合——在我们的延迟实验中，两者的组合表现优于纯重要性采样基准。